% LSPV-2 REPORT::COVERCOF1::2026-10-12T19:23:17Z
\documentclass[11pt]{article}
\usepackage[margin=1in]{geometry}
\usepackage{amsmath,amssymb,booktabs,array,hyperref}
\usepackage[T1]{fontenc}
\hypersetup{colorlinks=true,linkcolor=blue,urlcolor=blue,citecolor=blue}
\newcommand{\Status}[1]{\textbf{#1}}
\newcommand{\SAT}{\mathrm{SAT}}
\newcommand{\Dec}{\mathrm{Dec}}
\title{LSPV-2 COVERCOF1: A two-cofactor normal form falsifies universal post-coverage work}
\author{UNIVERSAL\_CORE\_TESTS (LSPV-2)}
\date{2026-10-12T19:23:17Z}

\begin{document}
\maketitle

\section*{Executive status}
\Status{PROVED/EXACT:} every total deterministic random-access decider has a
language-equivalent two-cofactor implementation that performs only \(O(1)\)
work after its first read of the final unread source bit. The transformation
costs two full calls to the original decider, \(O(N)\) input work, and linear
input storage. Consequently, the candidate universal lower bound that every
decider must perform superpolynomial work after complete real-input coverage
is \Status{FALSIFIED/RETIRED}.

This theorem does not reduce total SAT time, improve an exact-SAT base, or
establish a total-time or space lower bound. It is a quantifier-level
counterexample to one proposed solver-independent checkpoint.
\[
 \boxed{\text{P versus NP remains open.}}
\]

\section{Predecessor write-recovery preflight}
The newest canonical LSPV-1 artifact available to this run is
\path{research/agents/lspv-1/2026-10-13_0311_LSPV-1_HIERFIRE1.tex};
its required header and final sentinel verify. Repository history contains no
later canonical LSPV-1 report attributable to the immediately preceding
stagger. No exact missing scientific body was available from a canonical
artifact, so none was reconstructed. The gap is recorded for predecessor
recovery and did not suppress this scientific run.

\section{CORPUS REVIEWED}
The run surveyed the repository tree and read root \path{main.tex},
\path{AGENTS.md}, \path{CODEX.md}, and
\path{.codex/skills/pnp-research/SKILL.md}; all current specifications under
\path{tasks/}; actual \path{main4.tex} and \path{main5.tex}; the newest verified
MAIN snapshot \path{research/main/2026-10-13_0345_MAIN_REPPROMISE1.tex}; and
the newest verified LSPV-1 through LSPV-4 reports HIERFIRE1, SUCTYPE1,
REVUNCOMP1, and READFLOOR1. Repository search also covered materially related
historical reports TESTS-0014, PRODSTREAM1, DERIVEDLEAK1, TRACECONJ1,
READCEIL1 and the HEC/VSC-HEC/G-PPC-HEC/CRS-PPC/CCSPC/RCCN exploration
branch, together with governance notes and cited experiment artifacts. The
requested \path{tasks/TASK-008_SLICE-UNSAT.md} was not present and was not
reconstructed.

\section{Primary-literature and transfer preflight}
The two obstacles selected before substantive work were:
\begin{enumerate}
 \item transferring a semantic checkpoint into an obligation on every
 deterministic implementation; and
 \item matching the lower-bound model to arbitrary uniform polynomial-time
 random-access SAT deciders, including wrappers and alternate representations.
\end{enumerate}
Shannon cofactoring directly supports the Boolean identity used below, but not
a computational lower bound \cite{Shannon1949}. Decision-tree results such as
Rivest--Vuillemin apply exactly to their query model, not unrestricted RAM
computation \cite{RivestVuillemin1976}. Pippenger--Fischer provides an
oblivious simulation theorem with overhead, not an invariant internal
checkpoint \cite{PippengerFischer1979}. The Fortnow--Lipton--van
Melkebeek--Viglas SAT tradeoffs are model-specific time--space lower bounds
and do not supply the proposed universal suffix-work claim
\cite{FortnowEtAl2005}. Relativization and algebrization remain mandatory
barrier audits \cite{BGS1975,AW2009}.

\section{Exact model and falsified candidate}
Fix any decidable language \(L\subseteq\{0,1\}^*\). A decider is a total,
uniform, deterministic random-access machine \(M\) that outputs one bit and is
correct on every binary string. For a run on \(x\in\{0,1\}^N\), let
\(\tau_M(x)\) be the first step by which every real source position has been
queried, when such a step exists, and define
\[
 R_M(x)=T_M(x)-\tau_M(x).
\]
The candidate attacked in this run was:
\[
\begin{split}
\mathrm{POST\mbox{-}COVERAGE\mbox{-}SELECTION\mbox{-}WORK1}:\qquad
 &\exists\mathcal K\ \forall M\in\Dec(\mathcal K)\ \forall d\in\mathbb N\
 \forall N_0\\[-2mm]
 &\exists x\in\mathcal K,\quad |x|\ge N_0,\quad
 \tau_M(x)<\infty,\quad R_M(x)>|x|^d .
\end{split}
\]
Here \(\mathcal K\) may be an explicitly encoded clean-core family. The
universal quantifier over implementations makes the claim false for every
decidable \(\mathcal K\).

\section{DECISION-AT-COVERAGE-NORMAL-FORM1}
\paragraph{Theorem (\Status{PROVED/EXACT}).}
Let \(A\) be any total uniform deterministic random-access decider for \(L\),
running in time \(T_A(N)\) and work space \(S_A(N)\). There is a total uniform
deterministic decider \(B\) for the same language such that, for every
\(x\in\{0,1\}^N\), \(B\)'s first complete coverage of the real input occurs at
its last source read and only \(O(1)\) work remains. In standard word-RAM
accounting,
\[
 T_B(N)\le 2T_A(N)+O(N),\qquad
 S_B(N)\le S_A(N)+N+O(\log N).
\]
Polynomial time is preserved without fixing or hiding the polynomial exponent.

\paragraph{Construction.}
For \(N\ge1\), read and store \(x_1,\ldots,x_{N-1}\), deliberately withholding
\(x_N\). Simulate \(A\) on each of the two virtual strings
\[
 x^{(b)}=(x_1,\ldots,x_{N-1},b),\qquad b\in\{0,1\},
\]
answering virtual input queries from the stored prefix and hardwired bit
\(b\). Store the output bits \(a_0=A(x^{(0)})\) and
\(a_1=A(x^{(1)})\). Now read the real bit \(x_N\) for the first time and
output \(a_{x_N}\).

\paragraph{Proof.}
Exactly one virtual string equals the real input:
\(x^{(x_N)}=x\). Therefore \(B(x)=a_{x_N}=A(x)\). Before the final source
read only \(N-1\) real positions have been queried. That read completes
coverage, after which a two-entry table selection and one-bit output remain.
The two simulations and prefix copy give the displayed costs. Totality of
\(A\) ensures both virtual simulations halt, including on a syntactically
malformed cofactor. The \(N=0\) case is handled directly. \(\square\)

\paragraph{P=NP counterfactual and model matching.}
If \(\SAT\in\mathrm P\), apply the theorem to a polynomial-time SAT decider
\(A\). The resulting \(B\) is a model-matched uniform polynomial-time SAT
decider with constant post-coverage work. Hence no proof may derive an
unconditional contradiction merely from a universal lower bound on work after
complete source coverage; that alleged arrow already fails for every decidable
language. The theorem does not assume or imply \(\mathrm P\ne\mathrm{NP}\).

\section{Semantic compatibility and hostile controls}
For a CNF \(F\) and partial assignment \(\alpha\), define
\[
 H(F,\alpha)=\mathbf 1\!\left[\exists z\supseteq\alpha\; z\models F\right].
\]
The construction need not materialize \(H\), a witness, a separator, a
compatibility table over assignments, or a reconstruction object. It stores
only two terminal decision bits for two source cofactors. Thus large semantic
state count alone yields no suffix-time lower bound.

The identity survives COPY, DEAD, DELAY, interpreter-style wrappers,
recomputation wrappers, bit permutations, complement encodings, padding and
redundant clauses. It also applies to Horn, affine/Tseitin, 2-SAT,
bounded-treewidth, parity and independent-product controls. These controls do
not show general 3-SAT easy; they show that the implementation-level checkpoint
is not language invariant.

\section{Bounded reproducible experiment}
The script \path{experiments/review/decision_at_coverage_normal_form.py}
exhaustively tested sizes \(1\le N\le9\), every withheld coordinate, all
strings, nine predicate families, three encodings and six wrapper labels. It
performed \(1{,}548{,}666\) comparisons in \(1.1083\) seconds and found zero
failures. The output is
\path{results/raw/review/2026-10-13_decision_at_coverage_normal_form.json}.
The script SHA-256 is
\texttt{370f768a694602f97495a13ff19c177707c53a80a6ee0c94e65a14a29c437e39};
the result SHA-256 is
\texttt{66780d22d376b7656d90c4fb4a1a5744fabff57d2f0bcafaf698cbbfad562dfa}.
The invocation was designed to finish below 30 seconds; no chunking was
required. \Status{COMPUTATIONAL:} this is a finite regression check only. The
asymptotic theorem is established by the symbolic identity, not extrapolation.

\section{Faithful comparison with main4 and main5}
\path{main4.tex} develops exact CNF-level combinatorial quantities and guarded
structural statements. \path{main5.tex} develops RMV/numeric compatibility
machinery and explicitly records open acquisition and return arrows. Neither
document proves that an arbitrary SAT decider must expose a prescribed
semantic object or retain superpolynomial work after reading its final source
bit. COVERCOF1 does not reinterpret either manuscript; it removes one possible
but unsupported transfer principle.

\section{DUPLICATION / CONTRADICTION AUDIT}
Repository search found no prior statement of the exact two-cofactor
decision-at-coverage theorem. TESTS-0014 and PRODSTREAM1 already showed that
large product/semantic structure can coexist with a one-bit accumulator or
small streaming state. DERIVEDLEAK1, TRACECONJ1, READCEIL1 and READFLOOR1
study copying, schedules and source-query bounds. COVERCOF1 is distinct: it
targets the suffix after first complete source coverage and gives an
all-language compiler with fully charged costs. It is consistent with those
negative results and strengthens their warning that semantic size is not
runtime.

\section{WHAT CHANGED RELATIVE TO ROOT MAIN}
Before this run, root \path{main.tex} did not contain
DECISION-AT-COVERAGE-NORMAL-FORM1 or the retirement of
POST-COVERAGE-SELECTION-WORK1. The LSPV-2 region should now record the exact
two-cofactor compiler and all costs; the falsification of universal
post-coverage suffix-work lower bounds; the non-implication for total time,
memory, communication or SAT hardness; and the remaining need for a total-cost,
representation-robust transfer.

\section{ADOPTED / REJECTED / RETIRED PRIOR COMPONENTS}
\begin{center}
\begin{tabular}{>{\raggedright}p{0.24\linewidth}p{0.66\linewidth}}
\toprule
\Status{ADOPTED} & SAME\_REALIZATION, exact quantifiers, source accounting,
wrapper tests, arbitrary polynomial exponents, and the one-bit output audit.\\
\Status{ADOPTED} & Prior product-streaming and input-query negative controls as
supporting, not duplicate, evidence.\\
\Status{REJECTED} & Any inference from the number of semantic compatibility
states to time, space, communication, or exposed internal state.\\
\Status{FALSIFIED/RETIRED} & POST-COVERAGE-SELECTION-WORK1 and any equivalent
claim that every decider has superpolynomial suffix work after first complete
real-input coverage.\\
\bottomrule
\end{tabular}
\end{center}

\section{STALE OR OVERSTATED CLAIMS FOUND}
Any prose saying that a solver must still select among many semantic
possibilities after reading all source bits is overstated unless it names a
restricted model not closed under the two-cofactor compiler. The corrected
statement is only that total computation may remain hard; the location of that
work relative to source coverage is not invariant.

\section{Barrier-transfer and post-quantum ledger}
\begin{itemize}
 \item Shannon cofactoring: \Status{EXACT OBJECT/THEOREM TRANSFER}; it supplies
 the identity, not hardness.
 \item Pippenger--Fischer oblivious simulation and decision-tree query lower
 bounds: \Status{DIRECTLY TRANSFERABLE THEOREM} only within their stated
 models; use as unrestricted RAM lower bounds is a \Status{FAILED TRANSFER}.
 \item VLSI routing, communication cuts, pebbling and observability:
 \Status{ANALOGY/CANDIDATE}; a fully charged compiler is open.
 \item Schur squares, quadratic-relation ranks, Hilbert functions, syzygies and
 graded Betti data: \Status{FAILED TRANSFER} here. No theorem forces an
 arbitrary decider to compute these objects.
 \item Lossiness/collision entropy: the two-entry table has a trivial collision
 distribution and supplies no hardness. Generic collision entropy is a
 \Status{FAILED TRANSFER} absent a construction-specific extractor.
 \item Worst-case/average-case and collision-extraction lessons:
 \Status{ANALOGY/CANDIDATE} only; no unconditional black-box extraction theorem
 is obtained.
\end{itemize}
The construction relativizes and is purely Boolean, so it survives algebraic
extensions. It supplies no Natural-Proofs distinguisher. This is a
barrier-compatible falsification, not a route around the barriers.

\section{Strict status and highest-value next target}
\begin{center}
\begin{tabular}{lp{0.69\linewidth}}
\toprule
\Status{PROVED/EXACT} & two-cofactor normal form and cost accounting.\\
\Status{COMPUTATIONAL} & 1,548,666 bounded hostile-control checks, zero
failures.\\
\Status{FALSIFIED/RETIRED} & universal superpolynomial post-coverage work.\\
\Status{OPEN} & any superpolynomial total-time lower bound for an explicit
3-SAT clean core against every uniform deterministic decider.\\
\Status{OPEN} & an arbitrary-algorithm acquisition/return theorem invariant
under the wrapper and charging both full virtual runs.\\
\bottomrule
\end{tabular}
\end{center}
The single highest-value target is a total-cost lower-bound resource or
external extraction theorem invariant under equivalent wrappers and alternate
encodings. An exogenous-challenge HEC-style model remains a candidate only
after a fully charged compiler back to one static ordinary-CNF SAT decision.

\begin{thebibliography}{9}
\bibitem{Shannon1949}
C.~E. Shannon, The Synthesis of Two-Terminal Switching Circuits,
\emph{Bell System Technical Journal} 28(1), 1949.
\href{https://doi.org/10.1002/j.1538-7305.1949.tb03624.x}{doi:10.1002/j.1538-7305.1949.tb03624.x}.
\bibitem{RivestVuillemin1976}
R.~L. Rivest and J.~Vuillemin, On Recognizing Graph Properties from
Adjacency Matrices, \emph{Theoretical Computer Science} 3(3), 1976.
\href{https://doi.org/10.1016/0304-3975(76)90053-0}{doi:10.1016/0304-3975(76)90053-0}.
\bibitem{PippengerFischer1979}
N.~Pippenger and M.~J. Fischer, Relations Among Complexity Measures,
\emph{JACM} 26(2), 1979.
\href{https://doi.org/10.1145/322123.322138}{doi:10.1145/322123.322138}.
\bibitem{FortnowEtAl2005}
L.~Fortnow, R.~Lipton, D.~van Melkebeek, and A.~Viglas,
Time-Space Lower Bounds for Satisfiability, \emph{JACM} 52(6), 2005.
\href{https://doi.org/10.1145/1101821.1101822}{doi:10.1145/1101821.1101822}.
\bibitem{BGS1975}
T.~Baker, J.~Gill, and R.~Solovay, Relativizations of the P=?NP Question,
\emph{SIAM Journal on Computing} 4(4), 1975.
\href{https://doi.org/10.1137/0204037}{doi:10.1137/0204037}.
\bibitem{AW2009}
S.~Aaronson and A.~Wigderson, Algebrization: A New Barrier in Complexity
Theory, \emph{TOCT} 1(1), 2009.
\href{https://doi.org/10.1145/1490270.1490272}{doi:10.1145/1490270.1490272}.
\end{thebibliography}

% END_OF_VERIFIED_REPORT::LSPV-2::COVERCOF1
\end{document}
